\documentclass[13pt,oneside, aspectratio=169]{beamer}
\usepackage[utf8]{vietnam}
%\usepackage[linesnumbered, commentsnumbered]{algorithm2e}
\usepackage{appendixnumberbeamer}
\usepackage{mathptmx}

\usepackage{multirow}
\usepackage{makecell}
\usepackage{verbatim}


\usepackage[sorting=none,style=numeric,backend=bibtex]{biblatex}
\bibliography{reflib.bib}
%\usepackage{beamerthemesplit} 
\setbeamertemplate{footline}[frame number]
\setbeamertemplate{headline}{}
\usepackage{pgfplots}
\pgfplotsset{compat=newest}
\usepackage{booktabs}
%%%%%%%%%%%%%%%%%%
\definecolor{CTUcyan}{rgb}{0, 0.6862745098039216, 0.9372549019607843}
\definecolor{CTUblue}{rgb}{0.121568627450, 0.3607843137254902, 0.6627450980392157}

\usepackage{amsfonts,amsmath,amssymb}
\usepackage{algorithm}
\usepackage{algorithmicx}
\usepackage{verbatim}
\usepackage{graphicx}
\usepackage{subcaption}
\usepackage{booktabs,multicol,calligra}  %Allows the use of \toprule, \midrule and \bottomrule in tables
\usepackage{hyperref}
\usepackage{tcolorbox}
\usetikzlibrary{arrows,shapes}
\usepackage{multicol}
\usepackage{animate}



\mode<presentation>
{
	
	\setbeamertemplate{background canvas}{
		\begin{tikzpicture}
			\path (0,0) -- (\paperwidth,\paperheight);
			\node[opacity=.1] at (0.5\paperwidth,0.45\paperheight)
			{
\includegraphics[width=6cm]{BG/CTU_logo_singlecolor}};
		\end{tikzpicture}
	}
	
	\useinnertheme[shadow]{rounded}
	\usefonttheme{serif}
	\useinnertheme{circles}
	\setbeamertemplate{navigation symbols}{}
	\setbeamertemplate{caption}[numbered]
	\setbeamercolor*{palette tertiary}{use=structure,bg=CTUblue}
	\setbeamercolor*{palette quaternary}{fg=white,bg=CTUblue}
	\setbeamercolor{frametitle}{fg=CTUblue}
	\setbeamercolor{title}{fg=CTUblue}
	
	\setbeamercolor{block title}{bg=CTUblue, fg=white}
	\setbeamerfont*{frametitle}{size=\large,
		series=\bfseries,
		parent=structure}
	\setbeamercolor{block body}{bg=CTUblue!10}
	\setbeamercolor{section in toc}{fg=white}
	\setbeamercolor{subsection in toc}{fg=white}
	%	
}



\beamertemplatenavigationsymbolsempty




%\titlegraphic{
	%	\includegraphics[width=1cm]{Logo Đại Học Văn Lang H.png}
	%}
\title[]{\textbf{Nhận dạng thống kê cho hàm mật độ xác suất mất~cân~bằng\\ và ứng dụng cho ảnh}}
\subtitle{Báo cáo bảo vệ luận văn tốt nghiệp thạc sĩ}
\author[HUNG Tran-Nam]{Trần Nam Hưng\inst{1}}
\institute[CTU]{
	\inst{1} Khoa Khoa học Tự nhiên, Trường Đại học Cần Thơ
}

\date[Oct 26, 2025]{Ngày 26 tháng 10 năm 2025}

%\AtBeginSection[ ]
%{	
%	\usebackgroundtemplate{
\includegraphics[width=\paperwidth]{BG/thank2.pdf}}%
%			\begin{frame}
%						{\Large\bf\color{white} Mục lục}}
%						\tableofcontents[currentsection]
%					\end{frame}
%		}


\begin{document}
	
	{
		\usebackgroundtemplate{
\includegraphics[width=\paperwidth]{BG/background.pdf}}%
		\begin{frame}[plain]
			\vspace{2cm}
			\titlepage
		\end{frame}
	}
	
	{
		\usebackgroundtemplate{
\includegraphics[width=\paperwidth]{BG/thank2.pdf}}%
		
		
		\begin{frame}{}
			\vspace{3em}
{\Large\bf\color{white} Mục lục}
\large\color{white} \tableofcontents
		\end{frame}
	}

\section{Phần mở đầu}


\begin{frame}{1. Lý do chọn đề tài}
	
\end{frame}

\begin{frame}{2. Đối tượng và phạm vi}
	
\end{frame}


\begin{frame}{3. Phương pháp và cấu trúc}
	\begin{figure}
		\centering
		\includegraphics[width=\textwidth]{Figs/mindchart}
	\end{figure}
\end{frame}

\section{Phần nội dung}

\subsection{Vấn đề liên quan đến hàm mật độ xác suất}

\begin{frame}{1.1. Hàm mật độ xác suất}
	
\end{frame}

\begin{frame}{1.2. Hàm mật độ xác suất bổ sung}
	
\end{frame}

\begin{frame}{1.3. Độ đo khoảng cách hàm mật độ xác suất}
	
\end{frame}



\begin{frame}{1.4. Phân tích chùm hàm mật độ xác suất} 
	
\end{frame}


\begin{frame}{1.5. Phân loại hàm mật độ xác suất} 
	
\end{frame}





\subsection{Nhận dạng thống kê mất cân bằng}


\begin{frame}{2.1. Giới thiệu dữ liệu hàm mật độ}
	
\end{frame}

\begin{frame}{2.2. Phân tích chùm mất cân bằng (IFCF)}
	
\end{frame}

\begin{frame}{2.3. Phân tích chùm mất cân bằng (IFCF)}{Ví dụ số}
	
\end{frame}

\begin{frame}{2.4. Phân loại mất cân bằng (Bayes-IFCF)}
	
\end{frame}

\begin{frame}{2.5. Phân loại mất cân bằng (Bayes-IFCF)}{Ví dụ số}
	
\end{frame}



\subsection{Ứng dụng cho dữ liệu hình ảnh}


\begin{frame}{3.1. Phân loại mất cân bằng (Bayes-IFCF)}{Ví dụ hình ảnh}
	
\end{frame}

\begin{frame}{3.2. Phân loại mất cân bằng (Bayes-IFCF)}{Ví dụ hình ảnh}
	
\end{frame}

\section{Phần kết luận}

\begin{frame}{Kết luận và kiến nghị} 
	
\end{frame}




	\section{Tài liệu tham khảo}
	
	
	\begin{frame}[allowframebreaks, noframenumbering]{Tài liệu tham khảo}\small
		\printbibliography
	\end{frame}
	

	
	

{
	\usebackgroundtemplate{
\includegraphics[width=\paperwidth]{BG/thank2.pdf}}%
	\begin{frame}
		\begin{center}\color{white}
			\vspace{1.2cm}
			
			{\Large\color{white}\bf Chân thành cảm ơn quý Hội đồng!}
			
			
			\vspace{1cm}

		\end{center}
		
		
	\end{frame}
}



\end{document}